\section{引言：从 Reasoning 到 Agentic 的范式转移}

过去两年彻底改变了我们评估模型的方式以及对模型的期待。OpenAI 的 o1 证明了``思考''可以成为一种一等能力（first-class capability），即可以被专门训练并暴露给用户的能力。DeepSeek-R1 则证明了 reasoning 风格的后训练可以在原始实验室之外被复现和扩展。

\begin{importantbox}{核心论点}
2025 上半年的主旋律是 \textbf{reasoning thinking}——如何让模型花费更多推理时计算、如何用更强的奖励信号训练、如何暴露或控制额外的推理努力。而下一阶段的答案是 \textbf{agentic thinking}——\textit{thinking in order to act, while interacting with an environment, and continuously updating plans based on feedback from the world.}（为了行动而思考，在与环境交互的过程中，根据世界的反馈不断更新计划。）
\end{importantbox}

本文是 Qwen 团队 Junyang Lin 发表的一篇深度技术分析文章，从 o1/R1 的经验出发，反思了 reasoning 模型的局限性，并论证了 agentic thinking 作为下一代训练范式的必然性。文章视角独特，融合了 Qwen 团队的一线实战经验与对行业趋势的判断。


